\documentclass[a4paper,10pt]{article}
\usepackage[utf8]{inputenc}
\usepackage[T1]{fontenc}
\usepackage{lmodern}
\usepackage[margin=1.4cm, top=1.2cm, bottom=1.2cm]{geometry}
\usepackage{hyperref}
\usepackage{xcolor}
\usepackage{enumitem}
\usepackage{graphicx}
\usepackage{titlesec}

% Minimal clean styling
\hypersetup{colorlinks=true, urlcolor=blue!70!black, linkcolor=blue!70!black}
\pagestyle{empty}
\setlength{\parindent}{0pt}

% Section formatting
\titleformat{\section}{\large\bfseries\color{black}}{}{0em}{}[\titlerule]
\titlespacing*{\section}{0pt}{6pt}{4pt}

% Colors
\definecolor{accent}{RGB}{20, 50, 120}

% Role command
\newcommand{\role}[4]{%
  \vspace{2pt}\noindent
  \begin{minipage}[t]{0.70\textwidth}\raggedright\textbf{#1} \textbullet\ \textit{#3}\end{minipage}%
  \hfill
  \begin{minipage}[t]{0.28\textwidth}\raggedleft\footnotesize\color{gray}#2 \textbullet\ #4\end{minipage}%
  \vspace{1pt}%
}
\setlist[itemize]{leftmargin=1.2em, itemsep=1.5pt, topsep=2pt, parsep=0pt}

\begin{document}

% ================= HEADER =================
\begin{minipage}[c]{0.80\textwidth}
  {\Huge\bfseries Prabhavit Gupta}\\[3pt]
  {\color{accent}\large AI Engineer $\vert$ Desarrollo de Soluciones IA, Machine Learning, MLOps \& Python}\\[6pt]
  {\footnotesize +39 351 610 7807 ~$\vert$~
  \href{mailto:prabhavitg@gmail.com}{prabhavitg@gmail.com} ~$\vert$~
  \href{https://www.linkedin.com/in/prabhavit/}{linkedin.com/in/prabhavit} ~$\vert$~
  \href{https://github.com/prab-gupta}{github.com/prab-gupta} ~$\vert$~
  España / Venecia (Residencia Permanente UE $\vert$ Permiso Ilimitado en España $\vert$ 100\% Remoto)}
\end{minipage}\hfill

% ================= SUMMARY =================
\section{Summary}
Ingeniero de Inteligencia Artificial (AI Engineer) con sólida base en ingeniería de software y 3 años de experiencia comercial en el diseño, desarrollo e industrialización de aplicaciones basadas en algoritmos de IA, Machine Learning y modelos generativos. Especialista en la construcción de servicios backend robustos con Python y FastAPI, pipelines RAG con recuperación vectorial y despliegue en producción con Docker y prácticas MLOps. Acostumbrado a trabajar de forma autónoma resolviendo problemas complejos y colaborando estrechamente con equipos multidisciplinares. Residente permanente de la Unión Europea con autorización laboral ilimitada en España y disponibilidad inmediata para teletrabajo 100\%.

% ================= EXPERIENCE =================
\section{Experience}

\role{\href{https://flowspace.co}{FlowSpace.co}}{Mar 2024 - Aug 2026}{Lead AI Engineer \& Software Developer (Freelance)}{Londres, UK (remoto)}
\begin{itemize}
  \item Diseñó y desarrolló aplicaciones de Inteligencia Artificial en producción para más de 4.500 usuarios empresariales activos, integrando algoritmos de procesamiento de lenguaje natural y modelos generativos.
  \item Construyó pipelines de datos y búsqueda semántica vectorial (pgvector, Chroma) con técnicas de re-ranking para optimizar la relevancia y eliminar respuestas erróneas.
  \item Implementó prácticas de MLOps: suites de pruebas automatizadas, evaluación de modelos, control de versiones de datos y monitorización continua de latencias (<1.8s).
  \item Desarrolló microservicios y APIs REST escalables con Python y FastAPI, aplicando validación estricta de esquemas Pydantic y código limpio.
  \item Desplegó y gestionó contenedores Docker en infraestructura cloud (AWS / GCP) mediante flujos de integración y despliegue continuo (CI/CD).
  \item \textit{Tecnologías: Python, AI/ML, PyTorch, RAG, FastAPI, Docker, PostgreSQL, MLOps, CI/CD, Git}
\end{itemize}

\role{Urlink}{Apr 2025 - Jul 2025}{Software Engineer \& AI Developer (Freelance)}{Barcelona, España (remoto)}
\begin{itemize}
  \item Desarrolló algoritmos de procesamiento y clasificación automática de datos no estructurados para clientes internacionales.
  \item Construyó servicios backend asíncronos en FastAPI e implementó validaciones rigurosas para asegurar la calidad de los datos alimentados a los modelos analíticos.
  \item Colaboró con equipos técnicos de producto para definir especificaciones de integración y optimizar el rendimiento de las consultas.
  \item \textit{Tecnologías: Python, Scikit-learn, Pandas, PostgreSQL, Docker, REST APIs, Git}
\end{itemize}

\role{Consultor Técnico Independiente}{Sep 2023 - Mar 2024}{AI Solutions Engineer}{Venecia, Italia}
\begin{itemize}
  \item Desarrolló soluciones personalizadas de Inteligencia Artificial, automatización de procesos y asistentes conversacionales para pymes europeas.
  \item Analizó necesidades de negocio y diseñó arquitecturas técnicas viables y seguras para la adopción de tecnologías de IA.
  \item Redactó documentación técnica, especificaciones de APIs y manuales operativos para facilitar el mantenimiento de las soluciones implementadas.
  \item \textit{Tecnologías: Python, LangChain, OpenAI APIs, Vector DB, FastAPI, Linux, Git}
\end{itemize}

\role{\href{https://tractionboard.co}{Tractionboard.co}}{Aug 2023 - Oct 2023}{Full-Stack Automation Developer (Prácticas)}{Utrecht, Países Bajos (remoto)}
\begin{itemize}
  \item Desarrolló webhooks y flujos automatizados de sincronización de datos para plataformas SaaS.
  \item \textit{Tecnologías: Python, JavaScript, REST APIs, Git}
\end{itemize}

% ================= EDUCATION =================
\section{Education}
\role{Universidad Ca' Foscari de Venecia}{Oct 2020 - Jun 2024}{Grado en Filosofía, Estudios Internacionales y Economía (Minor en Emprendimiento)}{Venecia, Italia}

% ================= CERTIFICATIONS =================
\section{Certifications}
Microsoft: Introduction to building AI agents \textbullet\ Microsoft Copilot Studio Agent Academy \textbullet\ Microsoft Power Automate \textbullet\ Tray.AI Build Practitioner (Sep 2024) \textbullet\ Tray.AI Build Foundations

% ================= ADDITIONAL INFORMATION =================
\section{Additional Information}
\begin{itemize}
  \item \textbf{Autorización Laboral:} Residente Permanente de la Unión Europea con derecho laboral ilimitado en España (sin necesidad de visado).
  \item \textbf{Idiomas:} Inglés (C1 profesional fluido), Español (intermedio comunicativo), Italiano (profesional), Hindi (nativo).
  \item \textbf{Disponibilidad:} Inmediata (0 días de preaviso) para modalidad 100\% teletrabajo.
\end{itemize}

% ================= CORE SKILLS =================
\section{Core Skills}
\begin{itemize}
  \item \textbf{Inteligencia Artificial \& ML:} Algoritmos de IA, Machine Learning (PyTorch, Scikit-learn), LLMs, RAG, Búsqueda Vectorial, Prompt Engineering.
  \item \textbf{Ingeniería de Software \& Backend:} Python 3.12, FastAPI, Programación Asíncrona, Pydantic, Microservicios, APIs REST, SQL (PostgreSQL).
  \item \textbf{MLOps \& Despliegue:} Docker, CI/CD, Monitorización de Modelos, Evaluación Continua, AWS, GCP, Git, Linux.
  \item \textbf{Metodología:} Análisis de Problemas, Trabajo Autónomo, Buenas Prácticas de Código Limpio, Trabajo en Equipo Multidisciplinar.
\end{itemize}

\end{document}
